\documentclass[11pt,a4paper]{ctexart}
\usepackage[margin=2.4cm]{geometry}
\usepackage{amsmath,amssymb}
\usepackage{fancyhdr}
\usepackage{enumitem}
\linespread{1.18}
\setlength{\headheight}{15pt}
\setlength{\parskip}{0pt}

\pagestyle{fancy}
\fancyhf{}
\fancyhead[L]{\small Team 2207343}
\fancyhead[R]{\small \thepage}
\renewcommand{\headrulewidth}{0.4pt}
\fancypagestyle{plain}{\fancyhf{}\renewcommand{\headrulewidth}{0pt}}

\ctexset{
  section/format = \Large\bfseries\raggedright,
  subsection/format = \large\bfseries\raggedright,
  subsubsection/format = \normalsize\bfseries\raggedright
}
\setcounter{secnumdepth}{3}

\begin{document}

%====================== 摘要页 ======================
\thispagestyle{empty}
\begin{center}
{\small 2022 MCM/ICM\quad 选题 A\quad 队伍控制号 2207343\quad 摘要页（Summary Sheet）}\\[18pt]
{\LARGE\bfseries 最优功率分配——驶向未来}\\[8pt]
{\large （Optimal Power Allocation: Ride to The Future）}\\[10pt]
\end{center}

\begin{center}
{\large\bfseries 摘\quad 要}
\end{center}

\vspace{6pt}
谁会想到，东京奥运会公路自行车个人计时赛的冠军竟是一位数学家？信不信由你，数学确实做到了这一点。在本文中，我们将建立功率曲线的数学模型，以帮助车手赢得比赛。

在任务一中，我们基于生物学原理建立了一个功率—持续时间模型。该模型包含三个阶段：极限域（Extreme）、重度域（Severe）与高强度域（Heavy）。将检索得到的车手数据代入模型后，我们绘制了不同性别的冲刺型车手与计时赛车手的功率曲线。

在任务二中，我们首先利用任务一得到的功率曲线建立人体能量消耗模型，该模型刻画了车手在能量上限方面所受的约束。我们分析了自行车动力学模型，以确定骑行过程中输出功率与速度之间的关系 $P_{out}-v$。通过分析赛道特征，我们建立了一个分段非线性规划模型，其适用范围涵盖直道、弯道与坡道。对于大坡度路段或急剧弯道，我们逐一进行了细致的计算。最后，我们借助 Mathematica 求解出车手在任务一给出的三条赛道上骑行时的最优输出功率曲线，完赛时间如后所列。与冠军成绩相比，我们的模型滞后在 6\% 以内，体现了模型的准确性。

在自设计赛道上，男子冲刺型车手、女子冲刺型车手、男子计时赛车手与女子计时赛车手的完赛时间分别为 47.53 min、48.58 min、46.48 min 与 48.20 min；在日本赛道上，上述四类车手的完赛时间分别为 62.68 min、40.97 min、56.03 min 与 33.63 min；在比利时赛道上，四类车手的完赛时间分别为 58.15 min、59.62 min、53.03 min 与 54.05 min。

模型求解完成后，我们在任务三与任务四中进行灵敏度分析。对于天气因素，我们重点关注风速与风向的影响。当风速在 $\pm 2$ m/s 范围内变化、风向在 $\pm 180^{\circ}$ 范围内变化时，我们用 Python 得到输出功率随其变化的关系图。我们得出结论：模型对侧风的灵敏度在 1\% 以内，因而具有稳健性；而逆风对模型有约 10\% 的影响，这为团队比赛策略提供了思路。

我们还检验了模型对车手偏差的灵敏度。以比利时赛道为例，我们建立车手的实际输出功率函数 $P_{out}$，并加入随机扰动项 $g_{P_{out}}$ 以模拟车手的随机性。经计算，车手能够在体力耗尽之前完成比赛，时间滞后在 2\% 以内。

受模型对逆风高度灵敏这一特点的启发，我们在任务五中制定了团体计时赛策略。结果如下所列：均衡段（Balance，占 48\%）采用金字塔队形（Pyramid formation），储能段（Storage，占 27\%）采用一字长蛇阵列（Linear array），冲刺段（Sprint，占 25\%）采用盾形队形（Shield formation）。

最后，我们为车队经理（Directeur Sportif）撰写了一份面向计时赛车手与冲刺型车手的比赛指导。

\vspace{6pt}
\noindent\textbf{关键词：}功率曲线，非线性规划，动力学模型，队形规划

\newpage
\pagestyle{fancy}

%%====================== 正文 ======================
\section{引言}

\subsection{问题重述}

自行车运动的特点是在一段时间内保持高功率输出。这种能力可以用功率曲线来概括，它表示自行车手随时间所能产生的最大功率。由于人类耐力有限，功率曲线一般向右下方倾斜，如图 1 所示。

功率曲线用于为不同赛道制定配速策略，并且可以防止车手因一时冲动而承受显著的能量损失。这在任何需要配速的自行车比赛中都很有用，例如个人计时赛、Strava 赛段或比赛中的爬坡。在不同路况下使用不同的骑行功率，可以帮助车手为冲刺储备能量并赢得比赛。

图 1：功率曲线

我们正在为不同类型的自行车手建立一个数学模型。该模型受车手所能消耗的总能量限制，并聚焦于车手的若干特征，如骑行能力、车手姿势与车手的输出功率。基于上述理解，我们需要解决以下问题：

\begin{itemize}
\item 定义计时赛车手与冲刺型车手的功率曲线，并考虑性别因素。
\item 使用任务一中定义的每一条功率剖面，将模型应用于下列个人计时赛：
  \begin{itemize}
  \item 日本东京 2021 年奥运会个人计时赛赛道；
  \item 比利时佛兰德 2021 年国际自行车联盟（UCI）世界锦标赛个人计时赛赛道；
  \item 至少一条自行设计的赛道，其中包含至少四个急弯与不可忽略的道路坡度，且赛道终点应靠近其起点。
  \end{itemize}
\item 确定模型对天气因素（包括风速与风向）的灵敏度。
\item 分析结果对车手偏离目标功率分配的灵敏度。
\item 进一步讨论将模型推广到六人团体计时赛并得到最优结果。团队成绩由第四名车手越过终点线的时刻决定。
\item 为某车队的车队经理（Directeur Sportif）撰写一份两页的车手比赛指导，内容应聚焦于某条个人计时赛赛道与某位车手。它需要包含模型概述与方向性说明。读者是没有数学背景的人士。
\end{itemize}

\subsection{文献综述}

功率剖面研究一直是自行车运动研究的热点领域。自 20 世纪 80 年代末功率计发明以来，一些研究者选取了最能描述运动员训练强度范围的功率—持续时间模型。例如，双参数临界功率（CP）模型 [1] 只适用于重度运动强度域。然而，大量功率输出落在高强度运动强度域内，因此该模型并不能对车手的功率曲线作出公平的推广，这也说明不同的情形可能需要不同的建模方法。

随着生理学的发展，人们日益认识到功率曲线可以划分为三个不同的运动强度域：极限域、重度域与高强度域 [2]，这为研究可耐受的运动极限提供了综合性的方法。既有研究表明，无氧功率储备（anaerobic power reserve，APR）能够很好地预测极限运动强度域内的短期（$<3$ min）功率输出 [3]。在重度运动强度域内，主要使用 CP 模型，此外还有三参数 CP 模型 [4]。目前，针对高强度运动强度域已存在全功率模型（omni-domain power model）[5]。

\subsection{本文工作}

我们在假设与文献的基础上建立了两个模型，包括三阶段功率曲线模型与自行车动力学模型。功率曲线模型的三个阶段为极限域、重度域与高强度域，它们以人体不同的生理反应为特征。自行车动力学模型基于物理原理，描述天气因素对输出功率的影响。随后我们在不同的车手与赛道上检验了模型，结果良好。最后，我们设计了团体计时赛队形，使模型更具适用性。

图 2：本文的结构

\section{假设与合理性说明}

通过对问题的完整分析，为简化模型，我们作出以下合理假设。

\textbf{假设 1：}假设车手比赛场地的温度与湿度适宜，且波动范围较小。取温度范围为 $67.6\sim 75.6^{\circ}$F，湿度范围为 $35.1\sim 45.9\%$。

首先，考虑到包括题目给出的两项自行车赛事在内的大多数重大自行车比赛都在 7 月至 9 月举行；然后查阅以往天气状况，发现赛场的温度与湿度相对适宜。

\textbf{假设 2：}假设处于不同身体状态的车手以相同功率骑行相同时间，则消耗相同的能量。

专业运动员在比赛不同阶段体力分配的原理，大多基于乳酸堆积与 ATP 转化的模型 [6]。这些模型在解释人体能量消耗时具有线性性质 [7]。

\textbf{假设 3：}假设车手在骑行过程中恒定损失 5\% 的能量。

车手是提供功率的引擎。传动系统（链条、齿轮、轴承等）中的摩擦会造成少量损失，通常约为 5\% [8]；这里假设自行车的传动系统清洁且润滑良好。

\section{符号说明}

本文使用的主要数学符号列于表 1。

表 1：本文所使用的符号

\begin{itemize}
\item $CP$：临界功率，单位 W/kg；
\item $W'$：临界功率以上的做功能力，单位 J/kg；
\item $\lambda_E$：车手剩余能量与能量上限之比，无量纲；
\item $P_{out}$：车手的输出功率，单位 W/kg；
\item $t_n$：完赛时间，单位 s；
\item $v_{wind}$：风速，单位 m/s。
\end{itemize}

\section{三阶段功率—持续时间模型}

作为功率曲线计算与赛道专项决策的参数来源，功率剖面应当展现各类车手自行车比赛的重要生物学参数。基于这一前提，我们结合上述文献综述，分析现有模型的各种参数特征。最终，我们选取一名冲刺型车手与一名计时赛车手作为主要研究对象。我们还计算了 $CP$、$W'$ 等参数，它们有助于得到功率曲线。

接下来，我们获得不同车手在骑行 5 s、1 min、3 min、5 min 与 12 min 时所能维持的最大功率，其中一个数据源可用于参数计算。我们还查阅相关网站，获得最大有氧功率的耐受时间 $MAP_{TTF}$ 的计算方法，以及各阶段结束时的取值。

最后，我们计算了结果，并以表格形式呈现功率剖面，同时给出显示不同性别两名运动员功率曲线的图形。

\subsection{功率—持续时间模型的建立}

首先，我们用 $CP$ 表示车手的临界功率。临界功率是指车手在高强度下倾向于维持的功率输出，通常可持续 30 至 40 分钟。然后我们用 $W'$ 表示车手在能量完全耗尽之前能够在 $CP$ 以上消耗的功（单位 J/kg）。由于人体代谢复杂，且关于运动过程中相关生物力学与生化反应的研究仍有待突破，因此 $CP$ 与 $W'$ 的计算一般由实验数据获得。由文献 [9] 可知，$CP$ 与 $W'$ 的计算公式如下：

\[
CP=\frac{12\cdot P_{12}-3\cdot P_{3}}{9}
\]

\[
W'=0.24s\cdot\left(P_{3}-P_{12}\right)
\]

其中 $P_3$ 是车手全力骑行三分钟的平均功率除以其质量，而 $P_{12}$ 与 $P_3$ 的定义相同，但取 12 分钟。此外我们还有 $MAP_{TTF}$ 的公式

\[
MAP_{TTF}=\frac{W'}{P-CP}+\frac{W'}{P_{max}-CP}
\]

其中 $P_{max}$ 为极限功率。

根据相关资料 [9]，我们估计 $P_{max}\approx P_{5s}$。综合以上公式，我们得到表 2 所示的结果。

表 2：不同类型车手的功率剖面 I

\begin{itemize}
\item 冲刺型车手（男）：在 5 s、1 min、3 min、5 min、12 min 内的最大功率分别为 21.03、9.09、6.95、4.81、4.70 W/kg；
\item 冲刺型车手（女）：上述五个时长的最大功率分别为 16.40、7.39、6.17、4.17、4.13 W/kg；
\item 计时赛车手（男）：上述五个时长的最大功率分别为 16.89、8.74、7.14、5.53、5.46 W/kg；
\item 计时赛车手（女）：上述五个时长的最大功率分别为 13.39、7.11、5.85、4.83、4.75 W/kg。
\end{itemize}

表 3：不同类型车手的功率剖面 II

\begin{itemize}
\item 冲刺型车手（男）：功能阈值功率（FTP）为 3.91，$CP$ 为 3.94，$W'$ 为 0.54，$MAP_{TTF}$ 为 1.10，$CP_{TTF}$ 为 33.79；
\item 冲刺型车手（女）：FTP 为 3.39，$CP$ 为 3.45，$W'$ 为 0.49，$MAP_{TTF}$ 为 0.99，$CP_{TTF}$ 为 32.60；
\item 计时赛车手（男）：FTP 为 4.98，$CP$ 为 4.90，$W'$ 为 0.40，$MAP_{TTF}$ 为 0.97，$CP_{TTF}$ 为 40.29；
\item 计时赛车手（女）：FTP 为 4.38，$CP$ 为 4.38，$W'$ 为 0.26，$MAP_{TTF}$ 为 0.92，$CP_{TTF}$ 为 34.41。
\end{itemize}

根据不同的运动时间以及各类运动员在不同阶段的特征，我们参照文献综述求解三阶段的功率曲线。其适用范围与计算模型列举如下（见表 4）。

表 4：三阶段功率—持续时间模型

\begin{itemize}
\item 第一阶段：极限域，采用无氧功率储备模型
\[
P(t)=P_{3}+\left(P_{max}-P_{3}\right)\cdot e^{-kt},\qquad t<MAP_{TTF}
\]
\item 第二阶段：重度域，采用双参数临界功率模型
\[
P(t)=\frac{W'}{t}+CP,\qquad MAP_{TTF}<t<CP_{TTF}
\]
\item 第三阶段：高强度域，采用全功率—持续时间模型
\[
P(t)=\frac{W'}{t}\cdot\left(1-e^{t\cdot\frac{CP-P_{max}}{W'}}\right)+CP-A\cdot\ln\left(\frac{t}{CP_{TTF}}\right),\qquad t>CP_{TTF}
\]
其中结合曲线连续性，取 $k=0.7$、$A=0.1$。
\end{itemize}

\subsection{功率—持续时间模型的求解}

然后我们结合式（1）、式（2）以及表 2 中的数据，连同表 3 中的数据，将其代入表 4 中的公式。我们用 Mathematica 计算并绘制曲线，得到如图 3 所示的功率曲线。

图 3：不同类型车手的功率曲线。左：0～40 min 的功率曲线；中：0～1 min（极限域）的功率曲线；右：1～40 min（重度域与高强度域）的功率曲线

观察这些曲线可知，在同一项目中，男性车手的功率曲线始终位于女性车手功率曲线之上，这与男性车手体能更好的规律相一致。此外，在短时爆发（$t<1$ s）方面，冲刺型车手的功率显著高于计时赛车手；而在中长距离骑行比赛中，计时赛车手更占优势。

\section{个人计时赛的分段非线性规划}

\subsection{人体能量消耗模型的建立}

为了将任务一建立的功率曲线模型应用于赛道，我们接下来推导一个适合本问题背景的人体能量消耗模型。

设 $m\cdot P_{out}$ 为车手在 $t$ 时刻输出的功率，其中 $m$ 为车手体重，$P_{out}$ 为车手单位质量的输出功率。

图 4：任务二的图示。左：变功率输出的示意图；中：$P_{out}-t$ 的离散函数；右：$P_{out}-t$ 的连续函数

根据假设 2，我们按照所消耗总能量的比例进行计算。对于图 4 中间所示的功率离散取值，车手剩余能量与能量上限之比表示为

\[
\lambda_E=\frac{E_R}{E_m}=1-\sum_{i}\frac{\Delta t_i}{T_i}
\]

其中 $T_i$ 是身体能量充足时当前功率的最长保持时间。在图 4 中间的例子中，$\lambda_E$ 的结果为 $\lambda_E=1-\frac{10\ \mathrm{s}}{20\ \mathrm{s}}-\frac{5\ \mathrm{s}}{10\ \mathrm{s}}=0$。

将离散模型改为连续模型。对于图 4 右侧所示的连续 $P_{out}(t)$，切出一块宽度为 $dt$ 的近似矩形，矩形的高度为 $T(P_{out})$。接着用积分代替求和，得到 $t$ 时刻 $\lambda_E$ 的公式

\[
\lambda_E(t)=1-\int_0^{t}\frac{d\tau}{T\left(P_{out}(\tau)\right)}
\]

其中 $T(P)$ 是功率曲线表达式 $P(t)$ 的反函数，即 $T(P)=P^{-1}(P)$。

如果车手恰好在到达终点时耗尽能量，则存在如下约束：

\[
1-\lambda_E(t)=\int_0^{t}\frac{d\tau}{T\left(P_{out}(\tau)\right)}=1
\]

然而，在求解实际问题时，车手不可能严格遵循优化解所给出的详细方案。这可能导致车手在完赛之前就触及体能极限。因此，我们将等式右侧的 1 减小为 0.95，保留 5\% 的剩余体力以应对上述情形。

\subsection{输出功率与速度之间的数学关系}

我们对自行车进行动力学分析，受力如图 5 左侧所示，并建立如下公式：

\[
F_{total}=\left(F_{roll}+F_{slope}+F_{accel}+F_{wind}\right)/\eta
\]

其中 $F_{total}$ 为运动员的主动力，$F_{roll}$ 为滚动摩擦力，$F_{accel}$ 为加速产生的惯性力，$F_{wind}$ 为空气阻力，$F_{slope}$ 为坡度阻力（下坡时为负值）。根据假设 3，我们取传动效率 $\eta$ 为 95\%。

图 5：左：自行车受力分析；右：$P_{out}-v$ 关系图

根据物理学原理，各物理量的计算公式如下：

\[
F_{roll}=C_r\cdot Mg\qquad\left(0.0015<C_r<0.015,\ M=m+m_{bike}\approx 1.2\cdot m\right)
\]

\[
F_{slope}=Mg\cdot\sin\theta
\]

\[
F_{accel}=a\cdot M
\]

\[
F_{wind}=\frac{1}{2}C_W A\cdot\rho v_{wind}^{2}
\]

其中 $\rho$ 为空气密度。根据假设 1 中的温度与湿度取值，我们得到 $\rho=1.29\ \mathrm{kg/m^{3}}$。$C_W$ 为风阻系数，$A$ 为迎风面积，二者满足 $C_W A=0.39\ \mathrm{m^{2}}$，且该式仅对正面来流有效。

在无风条件下，

\[
v_{wind}=v_{bike}
\]

综合以上各式，再由 $m\cdot P_{out}=F_{total}\cdot v$，我们可以得到 $P_{out}$ 与 $v$ 之间的转换关系，如图 5 右侧所示。

\subsection{自设计赛道的比赛策略}

我们首先设计一条包含直道、下坡与急弯等基本路况的赛道，以获得在不同路况下应采用的策略。随后将该策略用于求解两场真实比赛的复杂赛道问题。

图 6：左：自设计赛道场景构建；右：自设计赛道的三维地图

我们使用 Webots 软件构建赛道场景，并绘制自设计赛道的平面图与三维地图，如图 6 所示。AB 段为一段短上坡，CD 段为长上坡，EF 段为长下坡。

\subsubsection{弯道路段的骑行策略分析}

首先，我们分析弯道骑行策略。我们希望车手以允许的最高速度通过弯道。这不仅能缩短通过弯道的时间，还能减少直道上的加速与减速时间。

设弯道内径为 $r$，路宽为 $d$，$R$ 为最优路径圆弧半径。根据向心力计算公式 $F=mv^{2}/r$，沿图中所给出的路线行驶可以获得最大的过弯速率。此时各物理量之间的关系如下：

积分后，我们得到过弯的最大速度

\[
v_{max}=\sqrt{\frac{\mu g\cdot\cos\alpha\cdot\left[d+r-r\cos\frac{\theta}{2}\right]}{1-\cos\frac{\theta}{2}}}\qquad(\ast)
\]

由于自设计赛道的弯道总长为 0.185 km，与直道总长 36.5 km 相比可以忽略。因此，我们在计算时间时先只考虑直道时间，忽略弯道中的速度变化问题。我们将车手在弯道前后的输出功率近似为阶跃变化。完成上述处理后，我们得到每个弯道的最优通过方案。

\subsubsection{直道路段的骑行策略分析}

接下来，我们分析直道骑行策略。我们假设有 $n_i$ 段水平直道（长度为 $L_1,L_2,\dots,L_{n_i}$）、$n_j$ 段下坡直道（长度为 $D_1,D_2,\dots,D_{n_j}$）以及 $n_k$ 段上坡直道（长度为 $U_1,U_2,\dots,U_{n_k}$）。水平直道的转换函数 $v-P_{out}$ 记为 $v^{(straight)}(P)$，倾角为 $\theta_j$ 的下坡直道对应的函数记为 $v^{(downhill)}_{\theta_j}(P)$，而 $v^{(uphill)}_{\theta_k}(P)$ 为倾角为 $\theta_k$ 的上坡直道对应的函数。于是任务转化为一个规划问题。

在满足能量约束的条件下求 $P_{out}(t)$，使完赛时间 $t$ 最小，即

\[
t=\sum_{i=1}^{n_i}\frac{L_i}{v^{(straight)}(P_i)}+\sum_{j=1}^{n_j}\frac{D_j}{v^{(downhill)}_{\theta_j}(P_j)}+\sum_{k=1}^{n_k}\frac{U_k}{v^{(uphill)}_{\theta_k}(P_k)}
\]

\[
\min\{t\}\qquad\mathrm{s.t.}\quad\int_0^{t_n}\frac{dt^{\ast}}{T\left(P_{out}(t^{\ast})\right)}=0.95
\]

在该映射中，$n_i=7$，$n_j=1$，$n_k=1$。$L$、$D$、$U$ 的下标从起点到终点按递增方式标注。

针对该函数，我们对 $P_{out}(t)$ 进行离散化处理。这使得 $\int_0^{t_n}\frac{dt^{\ast}}{T(P_{out}(t^{\ast}))}$ 的数值计算更易于实现。并且我们可以用 $P_i$、$P_j$、$P_k$ 绘制 $P_{out}(x)$ 图，这更有助于车手依据赛道景观判断输出功率。

图 7：坡道划分

图 8：计时赛男子车手的策略

我们使用 Mathematica 对反函数以及非线性规划函数 Minimize 进行数值求解，得到如表 5 所示的结果。

表 5：自设计个人计时赛赛道的求解结果

\begin{itemize}
\item 冲刺型车手（男）：$P_{out}(L)=6.0$，$P_{out}(U_1)=6.3$，$P_{out}(D_1)=5.2$，完赛时间 $t=47.53$ min；
\item 冲刺型车手（女）：$P_{out}(L)=5.8$，$P_{out}(U_1)=6.0$，$P_{out}(D_1)=5.1$，完赛时间 $t=48.58$ min；
\item 计时赛车手（男）：$P_{out}(L)=6.5$，$P_{out}(U_1)=6.4$，$P_{out}(D_1)=5.1$，完赛时间 $t=46.48$ min；
\item 计时赛车手（女）：$P_{out}(L)=6.0$，$P_{out}(U_1)=5.9$，$P_{out}(D_1)=5.0$，完赛时间 $t=48.20$ min。
\end{itemize}

我们接下来求解过弯时间。我们利用式（$\ast$）与前述转换函数进行求解，得到如下结果。

表 6：各弯道的求解结果

编号为 A、B、C、D、E 的五个弯道，其最优输出功率分别为 2.94、7.27、4.94、3.23 与 12.13 W/kg，对应的通过速度分别为 7.34、11.02、9.48、7.98 与 13.38 m/s。

比较上述两个表可以看出，由于弯道 A 与弯道 D 的角度较大，车手必须减速通过弯道。C 处的弯道虽然为 53 度，但它位于坡道上，静态摩擦力的上限降低，因此同样需要减速。弯道 B 则无需减速。

在自行车比赛中，加速与减速时的加速度可达 $2\ \mathrm{m/s^{2}}$，该过程所用时间可以忽略。

综合上述直道与弯道的结果，车手比赛过程中的最优 $P_{out}(t)$ 如图 8 所示（以男子计时赛车手为例）。
\subsection{2021 年 UCI 世界锦标赛个人计时赛赛道的策略}

图 9：2021 年 UCI 赛道分区

图 10：位于市中心广场附近的终点

根据 UCI 官方网站的数据，我们在 Google Earth 上确定了赛道所在区域的地形与路线特征。随后通过分析各路段的相似性与差异性，将赛道划分为 32 个部分（如图 9 所示）。其中，15 段为弯道，2 段为水平直道，3 段为下坡直道，12 段为上坡直道。

图 11：25 号路段的细节

值得注意的是，该赛区位于比利时西北部。它处于起伏较小的佛兰德高地平原。该区域上坡路段的最大坡度角为 $1.66^{\circ}$，而下坡路段的最大坡度角为 $0.69^{\circ}$。此外，相邻弯道之间有足够长的直道作为过渡。最短的直道是我们编号为 25 的路段（如图 11 所示）。25 号路段长约 400 m，足以让车手在较长时间内保持最优状态。因此，车手的通过时间几乎不受转弯的影响，分段离散规划模型仍然适用。

与东京奥运会个人计时赛赛道不同，本次比赛对选手的驾驶技术与过弯技术的要求相对较低，更注重体力分配与耐力。

\subsubsection{比利时赛道直道与坡道的骑行策略}

基于上述分析，我们仍可采用自设计赛道的方法。首先，我们求解直道策略：在满足能量约束的条件下求 $P_{out}(t)$，使完赛时间 $t$ 最小，即

\[
t=\sum_{i=1}^{n_i}\frac{L_i}{v^{(straight)}(P_i)}+\sum_{j=1}^{n_j}\frac{D_j}{v^{(downhill)}_{\theta_j}(P_j)}+\sum_{k=1}^{n_k}\frac{U_k}{v^{(uphill)}_{\theta_k}(P_k)}
\]

\[
\min\{t\}\qquad\mathrm{s.t.}\quad\int_0^{t_n}\frac{dt^{\ast}}{T\left(P_{out}(t^{\ast})\right)}=0.95
\]

其中 $n_i=2$，$n_j=3$，$n_k=12$。

考虑到上坡路段较多且坡度角各异，这会使模型难以收敛，对计算机的非线性规划求解极为不利。由于该区域地形起伏较小，我们选取了 9 段坡度角小于 $1^{\circ}$ 的路段（编号为 1、3、5、7、11、17、29、30、32），并将它们归入同一组。于是我们可以按路段长度对坡度角加权，得到综合而成的新坡度，随后用新坡度进行求解。

由于在小坡度角范围内 $\sin\theta$ 近似为线性，因此在合并多个坡道时，我们可以利用近似线性性质进行加权。相关参数由以下公式确定：

\[
U_C=\sum_{k\in A_U}U_k,\qquad \theta_C=\sum_{k\in A_U}w_k\cdot\theta_k
\]

\[
w_k=U_k\Big/\sum_{q\in A_U}U_q
\]

经过上述处理，问题转化为如下非线性规划问题：

在满足能量约束的条件下求 $P_{out}(t)$，使完赛时间 $t$ 最小，即

\[
t=\sum_{i=1}^{n_i}\frac{L_i}{v^{(straight)}(P_i)}+\sum_{j=1}^{n_j}\frac{D_j}{v^{(downhill)}_{\theta_j}(P_j)}+\sum_{k\notin A_U}\frac{U_k}{v^{(uphill)}_{\theta_k}(P_k)}+\frac{U_C}{v^{(uphill)}_{\theta_C}(P_C)}
\]

\[
\min\{t\}\qquad\mathrm{s.t.}\quad\int_0^{t_n}\frac{dt^{\ast}}{T\left(P_{out}(t^{\ast})\right)}=0.95
\]

图 12：2021 年 UCI 世界锦标赛个人计时赛赛道的求解结果。左：冲刺型车手的 $P_{out}(x)$ 图；中：计时赛车手的 $P_{out}(x)$ 图；右：弯道中最大速率的散点图

我们使用 Mathematica 求解每一路段的最优输出功率，运动员的曲线如图 12 的左图与中图所示。四类车手的完赛时间如表 7 所列。

表 7：2021 年 UCI 世界锦标赛个人计时赛赛道的求解结果

男子冲刺型车手的完赛时间为 58.15 min，女子冲刺型车手的完赛时间为 59.62 min，男子计时赛车手的完赛时间为 53.03 min，女子计时赛车手的完赛时间为 54.05 min。

\subsubsection{比利时赛道转弯处的策略}

此后，我们利用式（$\ast$）与 $P_{out}(v)$ 求解弯道中的最大功率，结果如图 12 的右图所示。图中的散点表示与 $x$ 相对应的弯道的最大转弯功率，折线则表示男子计时赛车手的 $P_{out}(t)$。由图可见，除需要轻微减速的 23 号与 32 号弯道之外，车手在直道上的功率均小于最大转弯功率。

在上述计算中，我们将 $P_{out}(x)$ 近似为离散函数。我们将在 6.2 节进一步分析选手的实际输出，如图 13 所示。这将证明我们在该问题上所采取的处理方法是合理的。

图 13：实际输出功率

图 14：2021 年奥运会赛道分区

\subsection{2021 年奥运会个人计时赛赛道的策略}

由图 14 可以看出，东京赛道位于富山市附近，该市三面被险峻的北阿尔卑斯立山连峰环绕。这里地形变化很大，弯道急促，与比利时赛道有明显不同。该赛事对车手的坡道技术提出了很高的要求。此外，男子车手需要绕赛道骑行两圈，而女子车手只需骑行一圈。

我们根据该区域的地形与坡道特征，将东京赛道划分为 68 个路段与 2 个区（大圈与小圈），划分情况如图 14 所示。其中，大圈中的大部分路段以及小圈中的部分路段仍可用先前的模型计算。对于另外两处路段（图 16 中 24—27 号的大爬坡曲线段与图 17 中 55—56 号的急弯段），我们将随后单独分析。

\subsubsection{日本赛道常规路段的策略}

根据我们对赛道的划分，在其余 62 个常规路段中，$n_i=20$，$n_j=25$，$n_k=17$。由于 $1-\cos 2^{\circ}<7\times 10^{-4}$，$\cos 2^{\circ}-\cos 4^{\circ}<2\times 10^{-3}$，因此我们对角度分别满足 $\theta<2^{\circ}$ 与 $2^{\circ}<\theta<4^{\circ}$ 的坡道应用 5.4 节中的加权简化方法。满足条件的路段编号集合为 $T_{U2}$、$T_{D2}$、$T_{U4}$ 与 $T_{D4}$，并且 $|T_{U2}|=8$，$|T_{D2}|=9$，$|T_{U4}|=9$，$|T_{D4}|=11$。最终，问题转化为如下非线性规划问题：

\[
t=\sum_{i=1}^{n_i}\frac{L_i}{v^{(straight)}(P_i)}+\sum_{j\notin T_D}\frac{D_j}{v^{(downhill)}_{\theta_j}(P_j)}+\sum_{k\notin T_U}\frac{U_k}{v^{(uphill)}_{\theta_k}(P_k)}+\sum_{T_U,T_D}\frac{U_{T_i}}{v^{(up/down)}_{\theta_{T_i}}(P_{T_i})}
\]

其中 $t=\Sigma_i+\Sigma_j+\Sigma_k+\Sigma_T$，$\Sigma_T$ 表示合并项。在满足能量约束的条件下求 $P_{out}(t)$，使完赛时间 $t$ 最小，即

\[
\min\{t\}\qquad\mathrm{s.t.}\quad\int_0^{t_n}\frac{dt^{\ast}}{T\left(P_{out}(t^{\ast})\right)}=0.95
\]

借助 Mathematica，我们得到 $P_{out}(x)$ 曲线，结果如图 15 所示。与 2021 年 UCI 赛道相比，东京赛道更为复杂，选手体力分配策略的难度也显著增加。

图 15：2021 年奥运会个人计时赛赛道的求解结果。左：男子车手的 $P_{out}(x)$ 图；右：女子车手的 $P_{out}(x)$ 图（其中不连续的部分为特殊路段，需要单独求解）

图 16：24—27 号路段的细节

\subsubsection{日本赛道特殊路段的策略}

\textbf{（1）24—27 号路段的大山弯}

这四个路段的参数如下所列：24 号路段长 0.1 km，转弯半径为 0 m，坡度角为 $6.65^{\circ}$，转弯角为 $0^{\circ}$；25 号路段长度约为 0 km，转弯半径为 9 m，坡度角为 $10.25^{\circ}$，转弯角为 $147^{\circ}$；26 号路段长 0.1 km，转弯半径为 0 m，坡度角为 $5.99^{\circ}$，转弯角为 $0^{\circ}$；27 号路段长度约为 0 km，转弯半径为 8.5 m，坡度角为 $4.22^{\circ}$，转弯角为 $79^{\circ}$。

此处的直道部分较为陡峭，例如 63 号路段的坡度角为 $4.4^{\circ}$。在此处车手的速度为 6.83 m/s，输出功率为 5.5 W/kg。我们将以此数据为基础制定直道策略。

接下来，我们分析转弯策略。该弯道是一处持续爬升的急弯，我们用式（$\ast$）进行计算。25 号路段的最佳过弯速度为 6.0 m/s，此时 $P_{max}(25)=6.32$ W/kg；27 号路段的最佳过弯速度为 6.4 m/s，此时 $P_{max}(27)=5.31$ W/kg。如此，该路段所耗时间为 31.2 s。

\textbf{（2）55—56 号路段的急弯}

这两个路段的参数如下所列：55 号路段长 0.2 km，转弯半径约为 0 m，坡度角为 $0^{\circ}$，转弯角为 $0^{\circ}$；56 号路段长度约为 0 km，转弯半径为 12 m，坡度角约为 $0^{\circ}$，转弯角为 $166.9^{\circ}$。

对于该弯道，其内径很小，转弯角高达 $166.9^{\circ}$。经计算，$v_{max}(56)=5.806$ m/s，$P_{max}(56)=4.32$ W/kg。在 55 号直道上，$P_{out}=5.6$ W/kg，$v=10.0$ m/s。我们计划在 56 号路段起点前 20 m 处开始减速，并在通过 56 号弯道之后立即加速回到原速度，该过程耗时 21.7 s。

图 17：55—56 号路段的细节

综合 5.5.1 与 5.5.2 两节，四名选手在优化策略下的结果如表 10 所列。

表 10：2021 年奥运会个人计时赛赛道的求解结果

男子冲刺型车手的完赛时间为 62.68 min，女子冲刺型车手的完赛时间为 56.03 min，男子计时赛车手的完赛时间为 40.97 min，女子计时赛车手的完赛时间为 33.63 min。

可以看出，显著起伏的路况增加了对车手爬坡能力与耐力的考验。冲刺型车手的表现明显低于计时赛车手，这一差距大于 UCI 赛道上的差距。

\section{灵敏度分析}

\subsection{天气的灵敏度分析}

天气条件包括温度、湿度、风速、风向等。

根据血液中肌酸激酶（CK）与血红蛋白（HB）的理论 [10]，在一定范围内，温度与湿度对运动员体能的影响很小。由假设 1 可知，这一结论可以被进一步证明是合理的。

其次，降雨对沥青赛道的滚动摩擦系数影响很小。降雨只是提高了对车手过弯技术的要求。文献 [11] 指出，降雨的负面影响可以通过训练在一定程度上被消除。

综上所述，我们对天气因素的灵敏度分析将聚焦于风速与风向。

我们将灵敏度分析分为直道路段与转弯路段两部分，以确定风向与风力可能带来的影响。道路坡度的影响主要体现在自行车的摩擦上，因此在灵敏度分析中我们统一假设道路坡度为 0\%。

\textbf{（1）直道路段}

在直道上，只有有效风速会影响骑行，而有效风速与车手的方向相同。因此，本部分我们只分析模型对风速的灵敏度。

根据自行车赛事对风速的要求，以及日本与比利时 7 月的平均风速，我们选取约 2 m/s 的风速进行分析。

我们用 Python 绘制了风速偏差 $\pm 1\%$ 范围内的最大输出功率曲线与风阻影响功率曲线，结果如图 18 所示。

图 18：风速的灵敏度

图 19：$P_{out}(v)$ 与 $P_W(v)$ 的变化

图 19 给出了风速从 1 m/s 变化到 4 m/s 时各曲线（由下至上）的变化情况。

由图可见，当车手在直道上匀速骑行时，偏差为 10\%。因此，模型在直道路段对风速较为敏感。

\textbf{（2）转弯路段}

接下来，我们分析模型在转弯过程中对风速与风向的灵敏度。我们用 Python 在风速与风向偏差 10\% 的范围内对模型进行扰动，结果如图 20 所示，其中 $\gamma$ 为局部风向角。

在转弯路段，风向变化对车手最大输出功率的影响量级为 0.1\%，而风速变化的影响量级为 0.01\%，二者均远小于直道路段。因此，模型在转弯路段对风速与风向的敏感度较低。

综上，风向与风速在一定范围内的变化不会影响模型的转弯策略。另一方面，模型对逆风更为敏感。在个人赛中，顺风时可以适当降低输出功率，逆风时则可以提高输出功率。

图 20：风速与风向的灵敏度。左：日本赛道，$v_{wind}=2$ m/s，$\gamma=-45^{\circ}$；右：比利时赛道，$v_{wind}=2$ m/s，$\gamma=180^{\circ}$

此外，这一现象为团体计时赛的策略提供了思路，我们将在第七节重点讨论。

\subsection{车手偏差的灵敏度分析}

在 5.3 至 5.5 节的 $P_{out}(x)$ 求解过程中，我们采用了理想化的近似离散值模型，因此车手在每个小路段上的输出功率保持不变。在实际情形中，由于车手操作的随机性，很难严格遵循最优功率曲线。本部分我们将检验车手偏离目标功率时模型的灵敏度。

首先，以 2021 年 UCI 的个人计时赛为例。我们依据各段的端点值对离散的 $P_{out}(x)$ 进行样条插值，并将其作为车手实际输出的函数 $P_{out}$。随后，我们加入随机扰动项

\[
g_{P_{out}}=\sum_{i=1}^{4}A_i\sin\left(0.1i\cdot t\right)
\]

其中 $A_i$ 由随机数确定，且 $0<A_i<5\%\cdot P_{out}^{(max)}$。于是车手的实际输出功率函数为 $P_{out}^{(real)}=P_{out}+g_{P_{out}}$。我们进行三次实验，得到三条 $P_{out}^{(real)}(x)$ 曲线，并将它们与理想的 $P_{out}(x)$ 曲线绘制在同一坐标系中，如图 21 的左图所示。

图 21：模型的求解结果。左：$P_{out}(x)$ 与 $P_{out}^{(real)}(x)$ 的曲线；中：不同策略下的 $1-\lambda_E$；右：不同策略下的比赛用时

我们将这四条曲线代入能量约束，结果表明所有车手都能完成比赛。就能否完赛而言，模型的稳定性极佳。具体结果如图 21 的中图所示。

接下来，我们计算了车手到达终点所需的时间（见图 21 的右图）。车手的结果均在理想输出曲线之下，但滞后时间处于最优时间的 2.5\% 以内。因此，比赛表现对车手实际策略的灵敏度较低。我们建议车手在每个路段的末端保持理想的功率曲线，这有助于减少滞后时间。

最后，我们将模型的计算结果与实际比赛冠军的成绩进行了比较，结果如表 11 所列。

表 11：锦标赛成绩与本文模型结果的比较

在男子公路自行车个人计时赛中，我们的模型结果为 56.03 min，冠军成绩为 55.08 min；在女子公路自行车个人计时赛中，我们的模型结果为 34.63 min，冠军成绩为 30.22 min。

模型结果略落后于世界锦标赛的成绩。这可能与我们用于确定功率曲线的实验数据样本来自平均水平车手而非顶尖车手有关。

与世界冠军成绩相比，我们模型的求解结果是合理的。这也表明，我们的模型具有使能力中等的车手接近冠军车手表现的效果，从而说明模型的有效性。

\section{团体计时赛的进一步讨论}

在 6.1 节中，我们得出结论：侧风对车手的影响很小，而逆风的影响较大。我们发现这一特征可应用于团体计时赛的优化。通过查阅文献，我们得到以下结论：

由于跟骑减阻（drafting）效应的存在，在一定距离（$<1$ m）之内，跟骑队列中第二名车手受到的风阻降至单人骑行的 60\%—70\%；而第三名车手的阻力降至约 50\%—63\%；沿车队向后，阻力将逐渐降至约 47\%—60\%。

此外，由于"亚音速上游扰动"的存在，在多车手编队的情形下，最后一名车手未必风阻最小。事实上，最后一名车手是跟骑队列中唯一未能从亚音速上游扰动中获益的车手。

我们将主要分析自行车队在直道上应对逆风的队形策略。我们计算了 1 km 直道上的团队成绩与个人成绩，以验证策略的合理性与有效性。

对于 6 人自行车队，我们有蛇形（snake）与金字塔形（pyramid）两种队形。对于 4 人冲刺自行车队，我们设计了如图 22 左图所示的盾形队形。热力图中的数值表示团队比赛中个人的风阻与个人赛同一状态下风阻之比 $K_f$。

图 22：左：团体计时赛中三种队形的示意图；右：个人计时赛中的团队策略

我们将比赛分为三个阶段来计算 20 km 直道上的团队成绩：均衡（Balance）、储能（Storage）与冲刺（Sprint）。各段长度分别为 $L_B$、$L_S$、$L_R$，且 $L_B+L_S+L_R=20$ km。

\begin{itemize}
\item 均衡段：6 名车手在 6 个位置上轮流领骑，从而均衡风阻造成的体力消耗。我们选择平均阻力系数 $K_f$ 最低的队形作为最优队形。计算得到蛇形队形 $K_f=67.17$，金字塔队形 $K_f=64.83$，因此我们选择金字塔队形。
\item 储能段：由于变换队形会使车手额外消耗体力，因此应尽可能保持队形。此外，冲线的四名车手应在阻力系数较小的位置上轮转，而另外两名车手（称为挡风者）应在阻力系数较大的位置上轮转。考虑到冲刺型车手在长蛇队形中的平衡，权衡之后我们选择长蛇队形。
\item 冲刺段：根据文献 [12] 的仿真结果，我们选择阻力系数最小的盾形队形。
\end{itemize}

接下来，以男子计时赛车手的数据为例进行计算。结果为：均衡段占 48\%，储能段占 27\%，冲刺段占 25\%，如图 22 所示。20 km 直道的个人赛用时 24.8 min，团体赛用时 21.2 min，团体成绩明显优于个人成绩。与 2018 年环法第三赛段冠军 BMC 车队 56.1 km/h 的平均速度相比，我们的模型在直道上的平均速度为 56.6 km/h，证明了模型的准确性。

图 23：左：侧风情形下的骑行策略；右：2016 年环土耳其赛中的斜线阵列

此外，对于侧风的情形，我们以右前方来风为例。此时队形应调整为斜线阵列，如图 23 所示。这有利于降低阻力系数并减小侧风扰动。

\section{模型评价与进一步讨论}

\subsection{优点}

\begin{itemize}
\item 求解高效：在求解最优策略时，我们充分利用了函数的线性或近似线性性质。在进行积分与离散近似时，求解速度与收敛速度都得到了大幅提升。
\item 结果准确：我们在模型简化时依靠数据与文献，并且对赛道的特殊部分进行了单独分析，这保证了精度。此外，模型的求解结果接近冠军成绩，验证了结果的准确性。
\item 普适性：我们采用的分段非线性规划模型适用于赛道信息与选手信息已知的比赛。除自行车个人计时赛之外，游泳、跑步等比赛也可以通过修改参数加以应用。
\item 稳定性与高容错性：我们对模型进行了灵敏度检验，验证了其稳定性。我们选取平均水平车手的参数，并将能量利用率设为 95\%。这些处理也提高了模型的容错性。
\end{itemize}

\subsection{缺点}

\begin{itemize}
\item 缺乏对抗因素：我们的模型主要分析无对抗的个人计时赛。建模过程缺乏对车手之间攻防策略的考虑。
\item 精度难以进一步提高：我们在模型中使用了数值方法与离散方法，这些方法限制了模型的精度。因此，我们在第六节与第七节中对模型作了进一步讨论以改善这一问题。结果表明，改进后的模型更为精确。
\end{itemize}

\subsection{结论}

本文对骑行策略与功率曲线之间的关系进行了细致分析。在介绍模型之前，我们分析了各类车手与不同种类赛道的特征，随后提出了三阶段功率曲线模型与自行车动力学模型。

接下来，为了利用所建立的数学模型优化自行车个人计时赛的策略，我们依据坡度、弯道等赛道特征，将复杂赛道划分为若干简单路段。然后利用非线性规划求解最优策略，我们对自制赛道、比利时赛道与东京赛道进行了分析。不同类型与性别的选手具有不同的策略。对于比利时赛道，模型求解结果与世界锦标赛成绩的偏差约为 8\%；对于东京赛道，模型结果的偏差约为 6\%。此后，我们分析了模型对天气因素与车手偏差的灵敏度。在天气因素中，风速与风向起主导作用。我们得出结论：模型对逆风更为敏感，这一结论有利于团体计时赛的队形分析。在研究对车手偏差的灵敏度时，我们用随机扰动项对模型进行扰动，得到三条功率曲线。车手在该扰动下完成比赛，用时差异很小，证明了模型的稳健性。

\section*{参考文献}

[1] Moritani T, Ata AN, Devries HA, Muro M (1981) Critical power as a measure of physical work capacity and anaerobic threshold. Ergon 24(5):339350.

[2] Burnley M, Jones AM (2007) Oxygen uptake kinetics as a determinant of sports performance. Eur J Sports Sci 7(2):6379.

[3] Weyand PG, Lin JE, Bundle MW (2006) Sprint Performance-duration relationships are set by the fractional duration of external force application. Am J Physiol 290(3):5865.

[4] Morton RH (1996) A 3-parameter critical power model. Ergon 39(4):611619.

[5] Puchowicz MJ, Baker J, Clarke DC (2020) Development and field validation of an omni-domain power-duration model. J Sports Sci 38(7):801813.

[6] Guan Zhongwei. Design and theoretical basis of lactic acid tolerance training for 400m secondary athletes. Journal of Shanghai Institute of Physical Education, 2001(S1):142-144.

[7] Wang Guiying. A study on the main physical fitness indicators and evaluation standards of Chinese men's 400-meter athletes. Journal of Ezhou University, 2007, 14(2): 69-72.

[8] Rainer Pivit. Drag Force Formulas. Radfahren 2/1990, pp. 44-46.

[9] Tom Bell. Could Critical Power be a better alternative to FTP for training? Published: July 15, 2021.

[10] Ji Jianmin, Cai Dapeng. Effects of temperature and humidity on the training effect of male canoeists. Journal of Shenyang Institute of Physical Education, 2011, 30(03): 68-71.

[11] Steve Westlake. Are You Faster When It Rains. 1 MAR 2016.

[12] Bert Blocken, Yasin Toparlar, Thijs van Druenen, Thomas Andrianne. Aerodynamic drag in cycling team time trials. Journal of Wind Engineering \& Industrial Aerodynamics 182 (2018) 128145.

\section*{附录：致车队经理的比赛指导}

亲爱的车队经理：

我们很荣幸向您展示模型的结果——一个名为 RIDERMAN 的系统。RIDERMAN 是一个数学优化系统，可以为车手的比赛给出有效的策略。"作为一名数学家，你习惯于独自解决问题，这也正是我对待自行车运动的方式。"这是 2021 年东京奥运会女子个人计时赛冠军安娜·范德布雷根（Anna van der Breggen）所说的话。她拥有剑桥大学数学博士学位，是一位数学家。在某种程度上，她的胜利说明了数学优化在体育竞赛中的重要性。

通过动力学建模、分段分析、非线性规划求解等步骤，可以根据赛道与选手的信息求解出一套最优输出功率策略系统。我们将其应用于东京奥运会赛道、比利时 UCI 赛道以及一条自制赛道，结果表明，即使是体能中等的运动员，使用我们模型给出的输出功率策略，也能达到接近世界锦标赛的成绩。

下面，我们将向车手介绍这套基于赛道信息与选手数据的系统 RIDERMEN，并根据我们的数学优化模型提出一些建议。

\subsection*{直道类型对照表（计时赛车手与全能型车手）}

\begin{itemize}
\item 坡度 $0\sim 2^{\circ}$：平地；
\item 坡度 $2^{\circ}\sim 4^{\circ}$：U4/D4 型；
\item 坡度 $4^{\circ}\sim 6^{\circ}$：U6/D6 型。
\end{itemize}

以上即我们的数学优化模型所给出的结论。

亲爱的车手，我是 RIDERMAN——你在路上的小助手。请先允许我问你一个问题。

如果你是计时赛车手或全能型车手（rouleur），你是计时赛车手还是冲刺型车手？

请看这里。

首先，针对某条特定赛道：让我们以弯道为节点把路段划分开，并按以下标准确定路段的类型（U 表示上坡，D 表示下坡）。

现在，让我们看一看 RIDERMEN 为应对不同道路类型所准备的策略。

\begin{enumerate}
\item 对于 90\% 以上路段为平地的赛道：在骑行过程中尽量均匀地分配能量是最佳方案。（RIDERMAN 已在下方为您绘制功率曲线与功率—速度转换曲线，用于确定您期望达成的实际且合适的平均输出功率。）此处给出功率曲线与功率—速度转换曲线。
\item 对于多坡道赛道：RIDERMAN 为每一类坡道准备了相应的策略。
  \begin{itemize}
  \item U4 型上坡：维持或略微提高通过原平地时的功率，即可在速度损失较小的条件下节省体力消耗；
  \item U6 型上坡：输出功率可暂时提高平地功率的 8\%—10\%，从而减少通过弯道前后的加速与减速时间；
  \item D4 型下坡：输出功率可降低约 24\%；
  \item D6 型下坡：输出功率可降低约 30\%。
  \end{itemize}
\end{enumerate}

RIDERMAN 已为您准备好独特的策略，助您登上巅峰。（提示：对于上述标签未描述的路段，您的计时赛经验将派上用场。）

现在轮到冲刺型车手登场了。RIDERMAN 专门为您准备了一套不同的策略组合，以便您更好地适应比赛并在计时赛中脱颖而出。

如果你是职业冲刺型车手，但心怀登顶计时赛的志向，请看这里。

首先，针对某条特定赛道：让我们像专家型车手那样划分路段，但要依据不同的标准（U 表示上坡，D 表示下坡）。

\subsection*{直道类型对照表（冲刺型车手）}

\begin{itemize}
\item 坡度 $0\sim 1^{\circ}$：平地；
\item 坡度 $1^{\circ}\sim 2^{\circ}$：U2/D2 型；
\item 坡度 $2^{\circ}\sim 4^{\circ}$：U4/D4 型；
\item 坡度 $4^{\circ}\sim 6^{\circ}$：U6/D6 型；
\item 坡度 $>6^{\circ}$：UM/DM 型。
\end{itemize}

\begin{enumerate}
\item 对于 85\% 以上路段为平地的赛道：在骑行过程中尽量均匀地分配能量是最佳方案。（RIDERMAN 还根据冲刺型车手的平均水平为您绘制功率曲线与功率—速度转换曲线，用于确定您期望达成的实际且合适的平均输出功率。）此处给出功率曲线与功率—速度转换曲线。
\item 对于多坡道赛道：RIDERMAN 同样为每一类坡道准备了不同的对应策略（利用冲刺型车手能够快速加速与减速的特点，轻松衔接不同状态以减少爬坡损失）。
  \begin{itemize}
  \item U2/U4 型上坡：提高输出功率（约 3\%）；
  \item U6 型上坡：输出功率提高平地功率的 5\%—8\%；
  \item UM 型：将速度降至约 60\%；
  \item D4/D6/DM 型下坡：输出功率分别降低约 24\%/30\%/50\%。
  \end{itemize}
\end{enumerate}

接下来，考虑到计时赛的过弯策略与冲刺的过弯策略差别较大，RIDERMAN 为您额外准备了一张表格与一个计算公式（式（$\ast$））。建议您填写弯道数据并计算最大过弯速度。

RIDERMAN 祝您在计时赛中有好运，别忘了根据实际情况灵活调整。

感谢您的阅读，您可以打印本文，也可以通过电子邮件发送给队员。您也可以在训练中重点关注这些指导，以提高训练的效率和目的性。

祝您的车队在即将到来的计时赛中好运。

\begin{flushright}
您诚挚的，\\
2207343 队
\end{flushright}

\subsection*{转弯统计工作表}

该工作表供车手填写各弯道的编号、内径（m）、坡度（度）与转弯角（度），并据此计算最大过弯速度。

\end{document}
